\section{Background and prerequisites}\label{sec2}

\subsection{Notational conventions}

All manifolds and fibre bundles are assumed to be smooth, Hausdorff, without boundary and connected, and all maps thereof assumed to be smooth. Given any manifold $M$, we use $\XF_{M}$ to denote the sheaf of smooth vector fields on $M$, and $C^{\infty}_{M}$ the sheaf of smooth, real-valued functions on $M$. If $\pi_{B}\colon B\rightarrow M$ is a fibre bundle, then we denote by $\Gamma_{\pi_{B}}$ the sheaf of sections of~$\pi_{B}$. Given a smooth map $f\colon M\rightarrow N$ of manifolds, and sheaves $\SS_{M}$ on $M$ and $\SS_{N}$ on $N$, we~denote by $f_{!}\SS_{M}$ the pushforward of $\SS_{M}$ and by $f^{!}\SS_{N}$ the pullback of $\SS_{N}$~\cite[p.~65]{hartshorne}. Given a (possibly time-dependent) vector field $X = X(t,-)$ on an open set $\OO$ in a manifold $M$, we use $\Fl^{X}$ to denote its \textit{flow}. That is, for $x\in\OO$ and $t_{0}\in\RB$, $t\mapsto\Fl^{X}_{t,t_{0}}(x)$ is the unique solution to the initial value problem
\begin{gather*}
\frac{\rm d}{{\rm d}s}\bigg|_{t}f(s;x) = X(t,f(t;x)),\qquad f(t_{0};x) = x
\end{gather*}
defined by the vector field $X$~\cite[p.~236]{leesmth}. For time-independent $X$, $\Fl^{X}_{t,t_{0}}$ depends only on $t-t_{0}$ and will be denoted simply by $\Fl^{X}_{t-t_{0}}$. Vector fields are assumed time-independent unless otherwise stated.

\subsection{Singular foliations}

We begin by recalling the standard sheaf-theoretic definition of a singular foliation.

\begin{Definition}\label{singdef}
Let $M$ be a smooth manifold. A~\textit{singular foliation} on $M$ is a subsheaf $\FF\subset \XF_{M}$ of $C^{\infty}_{M}$-modules on $M$ for which the following hold.
\begin{enumerate}\itemsep=0pt
\item $\FF$ is \textit{closed under Lie brackets} in the sense that for every open set $\OO$ of $M$, $\FF(\OO)$ is closed under the Lie bracket of vector fields.
\item $\FF$ is \textit{locally finitely generated} in the sense that for each $x\in M$, there exists an open neighbourhood $\OO_{x}$ of $x$ and a finite family $X_{1},\dots,X_{k}$ of elements of $\FF(\OO_{x})$ such that $\FF(\OO_{x})$ is the $C^{\infty}_{M}(\OO_{x})$-span of the $X_{i}$.
\end{enumerate}
\end{Definition}

One can alternatively describe a singular foliation of a manifold $M$ as a locally finitely generated submodule of the compactly supported vector fields on $M$ which is closed under Lie brackets, as in~\cite{iakovos2}. By~\cite[Remark 1.8]{GarZam} these definitions are equivalent. One of the most important facts regarding singular foliations is the Stefan--Sussmann integration theorem~\cite{stefan, sussmann}.

\begin{Theorem}%\label{stefansussmann}
Let $M$ be a manifold with a singular foliation $\FF$. Then $\FF$ integrates to give a~decom\-position of $M$ into smoothly immersed submanifolds called leaves.
\end{Theorem}
